\chapter{Hydrides and Small Molecules: Geometry First}\label{ch:hydrides}
\index{hydrides!alkali (Group 1)}\index{hydrides!alkaline earth (Group 2)}\index{hydrides!Group 14, tetrahedral}\index{hydrides!Group 16, water}\index{water molecule}\index{methane}\index{atomization energy}\index{LiH}\index{2hom pair}\index{2split pair}

Chapter~\ref{ch:h2} showed that the two electrons of H$_2$ can be carried in two parameter-free
constructions, 2split and 2hom, and that the exact energy lies between them. This chapter takes the
same constructions to molecules with a heavier atom. The criterion is \textbf{geometry}: bond lengths,
and the residual force on each nucleus at the observed geometry. Binding energies are reported as a
diagnostic of the construction, not as a target.

The results come from the article \emph{RealQM for Molecules: Validation in Progress}, which is kept up
to date as the validation proceeds; this chapter is a snapshot of it.

\section{An earlier version withdrawn}\label{sec:hydride-setup}

Earlier versions of this chapter reported atomisation energies of hydrides across groups 1, 2, 14 and
16 within 3--9\% of experiment, with a single softening radius $r_c$ swept down each column. That
result is withdrawn, for four reasons found on rechecking it.
\begin{itemize}
\item The quantity was $E(R) - E(2R)$, not an atomisation energy. At $2R$ the model molecule is still
  strongly bound: rerun for NaH on the original page, the difference is about 60~kcal/mol, while the true
  binding $E(\mathrm{Na}) + E(\mathrm{H}) - E(\mathrm{NaH})$ is $0.1745\,E_h = 109.5$~kcal/mol --- the
  binding of H$_2$, not of NaH.
\item The bond lengths were not the experimental ones. The runs were made at $R = 1.0$ and $2.0$~\AA{} for
  every $r_c$; the NaH bond is $1.887$~\AA.
\item $r_c$ was chosen per molecule to match the target. A radius chosen that way is a fit, and the
  ``periodic trend'' in $r_c$ was the order of the fits. For Na the chosen $0.5$ lies below the radius of
  its own core charge.
\item Some constructions put more than two electrons in one three-dimensional domain (the $+2$
  carbon with its valence spread over four lobes), which the model does not allow (below).
\end{itemize}
Model NaH also has its dipole of the wrong sign. What replaces these tables is smaller, and every number
in it can be rerun.

\section{The rules}\label{sec:hydride-rules}

\paragraph{A domain holds one electron or a pair.} RealQM carries no spin variable; spin enters only
through the Pauli principle, as the rule that a three-dimensional domain holds at most two electrons, and
two only if of opposite spin. A pair can therefore be carried in two ways. In \emph{2split} each electron
owns its own domain and the two meet at a free wall: complete avoidance. In \emph{2hom} both share one
domain and one density: no avoidance beyond the mean field, which is the restricted Hartree--Fock pair. A
pair that shares one region --- a core $1s^2$, a lone pair, a bonding pair near equilibrium --- is 2hom.
Electrons that must avoid each other --- H$^-$, a bond pulled apart --- need 2split. The two are distinct
states with no parameter joining them, and both are reported.

\paragraph{Inner shells are excluded by a fixed radius.} An inner shell cannot be carried as a domain with
a free wall: the free wall charges nothing for crowding, the valence presses into the core and the bond
collapses. Every all-electron LiH with free core walls does this. The core is instead excluded by a fixed
sphere of radius $r_c$ around the nucleus, which then acts as a pseudo-nucleus of charge equal to the
valence count. The radius is set once per atom, from the atom: the radius at which the pseudo-atom alone
reproduces its ionisation energy. Nothing is adjusted to molecular data. This is a frozen-core model,
stated as such.

\paragraph{Energies are checked independently.} Above five domains the solver's default
electron--electron path was found to under-count the repulsion by several hartree (methane: about
$2.5\,E_h$), so every energy above five domains is recomputed from the converged density by an
independent program before it is quoted. The results of this chapter that involve more than five domains
were rerun on the per-electron path and checked that way.

\section{Light hydrides}\label{sec:group1}\label{sec:group2}

\begin{table}[h]
\centering\small
\setlength{\tabcolsep}{4pt}
\caption{Bond lengths ($a_0$) and binding energies ($E_h$, diagnostic) of light hydrides, three-dimensional,
$h = 0.1$. Bond length from the 2hom construction.}\label{tab:hydride-summary}
\begin{tabular}{llrrrrrr}
\toprule
molecule & construction & $R_e$ & exp. & error & $D_e$ 2hom & $D_e$ 2split & exp. \\
\midrule
H$_2$   & $1{+}1$ (axisym.)            & $1.368$ & $1.401$ & $-2.3\%$ & $0.138$ & $0.281$ & $0.1745$ \\
LiH     & Li$^+$ ($r_c$) + pair         & $2.85$  & $3.015$ & $-5\%$   & $0.095$ & $0.063$--$0.073$ & $0.092$ \\
BH      & B $2{+}3$ ($2_{\rm hom}{+}1$) & $2.29$  & $2.329$ & $-1.6\%$ & $0.155$ & $0.276$ & $0.135$ \\
OH      & O $+3$ ($2_{\rm hom}{+}1$)    & $1.79$  & $1.832$ & $-2\%$   & $0.241$ & $0.371$ & $0.170$ \\
BeH$_2$ & Be $2{+}2_{\rm split}$        & $2.44$  & $2.52$  & $-3\%$   & $0.248$ & n.c.    & $0.2231$ \\
\bottomrule
\end{tabular}
\end{table}

The bonds are $2$--$5\%$ short. The binding lies between the two constructions or near the 2hom value; the
2hom bonding pair is the closer for LiH and BH. (BeH$_2$: atomisation of both bonds; n.c., the 2split run
does not converge.)

\paragraph{LiH.} LiH is the simplest polar bond with an inner shell. Pseudo-Li carries charge $+1$ outside
a fixed core of radius $r_c = 1.99\,a_0$, the radius at which the pseudo-atom reproduces the Li ionisation
energy; the one-dimensional all-electron Li atom places its own core--valence interface at
$1.98$--$2.05\,a_0$, within 3\% of it. The two valence electrons form one 2hom pair over pseudo-Li and a
bare proton. The energy curve gives $R_e = 2.85\,a_0$ ($-5\%$) and $D_e = 0.0947\,E_h$ ($+3\%$), and at
$R = 3.0$ the dipole is $6.53$~D against $5.88$: of the right sign, Li$^+$H$^-$, and $11\%$ high. Nothing
was adjusted to LiH. With one electron per domain instead, the Li valence electron settles \emph{behind}
Li and the dipole has the wrong sign, whatever the starting density.

\paragraph{Forces at the observed geometry.} The residual force on each H along its bond at the observed
geometry is small and inward: $-0.022$ (LiH), $-0.013$ (BH), $-0.018$ (BeH$_2$), $-0.028\,E_h/a_0$ (OH).
The model's X--H bonds are slightly short, which is what the bond lengths say.

\section{Linear molecules}\label{sec:group14}

The same pseudo-atoms, unchanged, build linear molecules. In the axisymmetric reduction, HCN, CO and
CO$_2$ have their bonds within 3\% of experiment. Released as a whole from the observed geometry,
cyanoacetylene HC$\equiv$C--C$\equiv$N keeps C--H and C$\equiv$N within $1.3\%$, but the two
carbon--carbon bonds come out nearly equal: each bond is carried by one pair whatever its order, so a
triple and a single C--C bond are the same construction. Bond alternation is not represented, and the
C$\equiv$C of C$_2$H$_2$ comes out $10\%$ long for the same reason.

\section{Heavy atoms with their whole valence}\label{sec:group16}\label{sec:rc-meaning}

Carbon, nitrogen and oxygen carrying their whole valence (CH$_4$, NH$_3$, H$_2$O, ethane) do not yet work:
\begin{itemize}
\item The X--H bonds are a few per cent short (CH$_4$: $-4\%$), consistent with the light hydrides.
\item The binding is two to three times too large, by an amount per heavy atom: the excess of ethane is
  twice that of methane.
\item Shapes are not selected: the CH$_4$ tetrahedron and the square plane have the same energy within
  $0.002\,E_h$, and the H$_2$O angle depends on how the lone pairs are carried.
\end{itemize}
An earlier account blamed this over-binding on a missing Pauli kinetic energy of strong cores. That
account rested on the faulty electron--electron path and is withdrawn; with the repulsion counted
correctly the over-binding is smaller but remains. O$_2$ and F$_2$ are not reproduced by any core
configuration tried.

Most of these failures share one feature: where pairs are pressed together the free wall between them
charges nothing, and the model has no repulsion between pairs beyond their Coulomb energy. Whether RealQM
needs such a repulsion, and in what form, is the main open question for molecules. Its test is F$_2$
weakly bound near $2.67\,a_0$, the H$_2$O angle and the CH$_4$ tetrahedron recovered, and LiH, Li$_2$,
C$_2$, N$_2$ left where they are.

\section{What is established}\label{sec:hydride-summary-text}

\paragraph{A geometry model.} On light hydrides, homonuclear X$_2$ from Li$_2$ to N$_2$ and linear
molecules built from the same pseudo-atoms, bond lengths are within a few per cent, with small inward
forces at the observed geometry. Nothing is adjusted to molecular data; each atom contributes one core
radius, set from its ionisation energy.

\paragraph{Energies as a diagnostic.} Binding energies lie between, or near one of, the two constructions
of a pair. They are not chemically accurate, and above five domains they are quoted only after an
independent recomputation.

\paragraph{Open.} Heavy atoms with their whole valence, shapes and angles, O$_2$ and F$_2$, bond order,
dipoles beyond LiH, and reaction energies. These are the next steps, not results.

The next chapter goes to hydrogen-bonded dimers and reactions, where the model is used beyond what this
chapter validates.

% --- End of Chapter 8 ---
